\documentclass[10pt,a4paper]{amsart}
\usepackage[margin=19mm]{geometry}
\usepackage{amsmath,amssymb,mathtools}
\usepackage[scr=rsfs,cal=euler]{mathalfa}
\usepackage{xfrac}
\usepackage[unicode,hidelinks,bookmarks=false]{hyperref}
\newcommand{\ie}{i.e.\ }
\newcommand{\cf}{cf.\ }
\newcommand{\resp}{respectively\ }
\newcommand{\Subsection}{\subsection}
\providecommand{\nobreakdash}{\nobreak\hbox{-}}
\setlength{\emergencystretch}{2em}
\multlinegap0pt
\usepackage{amssymb}
\newtheorem{prop}{Proposition}
\newcommand{\C}{\mathbb{C}}
\newcommand{\Q}{\mathbb{Q}}
\newcommand{\Z}{\mathbb{Z}}
\title{Applications of representation theory and of explicit units to Leopoldt's conjecture}
\author{Fabio Ferri, Henri Johnston}
\date{Source: arXiv:2301.05700v2 (CC BY 4.0)}
\begin{document}
\maketitle

\noindent\emph{Source and licence.} Applications of representation theory and of explicit units to Leopoldt's conjecture. Version arXiv:2301.05700v2 (CC BY 4.0), distributed by arXiv under the Creative Commons Attribution 4.0 International licence, http://creativecommons.org/licenses/by/4.0/, as recorded in the arXiv licence metadata for this exact version and shown on the abstract page https://arxiv.org/abs/2301.05700v2. Full licence notice: \texttt{licenses/arxiv-2301.05700-CC-BY-4.0.txt}.\par\smallskip

\noindent\textit{Editorial scope.} Counted unit: the article's prop prop:infinitepol (source0.tex lines 1797--1803). The article's complete original proof of it is delivered with it (source0.tex lines 1805--1843, one \texttt{proof} environment of 39 lines). No mathematics is written, repaired or re-derived: the statement and the proof are the article's own lines, in the article's own notation, and no retained line is substituted or re-flowed.  No further source-local context is needed: every cross-reference of every retained line resolves inside the two delivered spans.  Nothing was omitted from any retained span, so the package carries no omission record.  No retained line contains a citation, so the wrapper carries no bibliography and no bibliography entry.  No retained line contains a figure environment, an image-inclusion command or drawing code, so no image is delivered and the package declares no asset dependency.  Editorial formatting only, in addition: an AMS \texttt{amsart} preamble that restates the article's own declarations and macros used by the retained lines, namely the environments \texttt{prop} on separate counters, together with the standard packages those lines need.  The retained environments print in this document as Proposition 1 in order of appearance, whereas the article prints the same items with the numbering of its own full document.  The complete native source of the article as distributed in the arXiv e-print for this exact version is bundled unchanged as \texttt{sources/136072/source0.tex}.
\begin{prop}\label{prop:infinitepol}
Let $g(y)\in\Z[y]$ be a polynomial that is not a square in $\C[y]$. 
Let $d \in \Z$ and let $I=\Z_{\geq d}$ or $\Z_{\leq d}$.
Let $a,b \in \Z$ with $a \neq 0$ and let $J=\{ ax+b \}_{x \in \Z} \cap I$ be an arithmetic progression in $I$.
Then there exists an infinite set of prime numbers $\mathcal L$ such that for each $\ell \in\mathcal L$, 
there exists $k \in J$ such that $\ell^{m} \parallel g(k)$ for some odd $m \in \Z_{>0}$.
\end{prop}

\begin{proof}
By the hypothesis and Gauss's lemma, we can write $g(y)=w(y)^{2}u(y)$ where $w(y),u(y) \in \Z[y]$ and
$u(y)$ is non-constant and square-free in $\C[y]$.
Set $v(x)=u(ax+b)$, which is a non-constant polynomial in $\Z[x]$. 
The values of $x$ such that $ax+b\in I$ are in a ray $R$ of the form $\Z_{\geq e}$ or $\Z_{\leq e}$ for some
$e \in \Z$.
It suffices to show that there exists an infinite set of prime numbers $\mathcal{L}$ such that for each
$\ell \in \mathcal{L}$, there exists $k \in R$ such that $\ell \parallel v(k)$.
After a straightforward linear reparameterisation, we can and do assume that $R=\Z_{\geq 0}$. 
Note that $v(x)$ has no repeated roots in $\C$, since it is a linear reparameterisation of $u(y)$.
Hence there exist $a(x),b(x)\in \Q[x]$ such that
\[
 a(x)v(x)+b(x)v'(x)=1,
\]
where $v'(x)$ denotes the formal derivative of $v(x)$. 
By clearing denominators, this implies that there exist $A(x),B(x) \in \Z[x]$ and $C \in \Z$ with $C \neq 0$
such that 
\begin{equation}\label{eq:lin-comb-formal-deriv}
A(x)v(x)+B(x)v'(x)=C. 
\end{equation}

It is well-known that if $h(y) \in \Z[y]$ is non-constant, 
then the set of prime numbers that divide some element of $h(\Z_{\geq 0})$ is infinite 
(sketch: consider $h(n!c^2)$ for sufficiently large $n$, where $c$ is the constant term of $h$).
In particular, this result applies to $v(x) \in \Z[x]$.
Now let $\ell$ be a prime number such that $\ell \nmid C$ and there exists $k_{0} \in \Z_{\geq 0}$ with 
$\ell\mid v(k_{0})$. 
We claim that we can always find $k_{1} \in \Z_{\geq 0}$ such that $\ell \parallel v(k_{1})$. 
Since there are infinitely many choices for $\ell$, this will prove the desired result.
If $\ell \parallel v(k_{0})$ then we take $k_{1}=k_{0}$.
Suppose that $\ell^{2} \mid v(k_{0})$. Then
\[
 v(k_{0}+\ell)-v(k_{0})=v'(k_{0})\ell+r\ell^{2}
\]
for some $r \in \Z$.
However, since $\ell\mid v(k_{0})$ and $\ell\nmid C$, we have that $\ell\nmid v'(k_{0})$ by 
\eqref{eq:lin-comb-formal-deriv}. 
Therefore $\ell^{2} \nmid v(k_{0}+\ell)$, and so we can take $k_{1}=k_{0}+\ell$.
\end{proof}

\end{document}
